% PDF p. 272
\section{引言}

流形假设（manifold hypothesis）是机器学习中广为接受的一项原则。它认为，表面上维数很高的数据其实集中在嵌入高维空间的低维流形附近（Whiteley et al.\ 2025）。流形假设指出，实际应用中遇到的高维数据常位于这类低维流形附近。这一假设对机器学习和统计推断具有深远影响，尤其是在传统欧氏方法不足以处理数据的情形。本章介绍流形上统计理论的最新进展、它们在人工智能中的应用、常用的流形假设检验方法、渐近理论和抽样技术，并讨论现有局限与未来研究所面临的挑战。

\subsection{流形上的统计理论}

信息几何将微分几何应用于概率论与统计学，着重研究统计流形（statistical manifold），其中每一点代表一个概率分布。这套框架提供了理解统计模型几何结构的工具，使经典统计方法能够更自然地推广到非欧氏情形。

\subsubsection{流形上的假设检验}

流形数据的几何结构是非欧氏的，因此在流形上检验假设需要专门的统计方法。
\begin{itemize}
  \item \textbf{内在方法（intrinsic method）}利用流形的黎曼结构，借助弗雷歇均值、测地距离等工具进行统计检验。
  \item \textbf{外在方法（extrinsic method）}将流形嵌入更高维的欧氏空间，采用传统的欧氏统计技术，再将结果投影回流形。
\end{itemize}
选择内在方法还是外在方法，取决于计算可行性、流形的具体几何性质等因素。

\subsection{流形上统计量的渐近理论}

流形上的渐近分析把经典统计结论推广到非欧氏情形。
% PDF p. 273
\begin{itemize}
  \item \textbf{相合性与正态性}在一定条件下，弗雷歇均值等估计量具有相合性和渐近正态性，这与欧氏情形类似。
  \item \textbf{曲率的影响}流形的曲率可能影响估计量的渐近分布，统计推断过程需要相应调整。
\end{itemize}

\subsection{流形上的抽样方法}

在流形上抽样对于概率建模和模拟等任务至关重要。
\begin{itemize}
  \item \textbf{流形上的 MCMC}改造马尔可夫链蒙特卡罗（Markov chain Monte Carlo，MCMC）方法，使其符合流形的几何性质，从而生成有效样本。
  \item \textbf{测地游走（geodesic walk）}沿流形上的测地线模拟随机游走，在生成样本时保留内在性质。
  \item \textbf{谱方法（spectral method）}利用流形的谱性质设计抽样算法，使样本能够体现其结构。
\end{itemize}

\subsection{局限与未来的研究挑战}

尽管已有进展，仍有若干问题有待解决。
\begin{itemize}
  \item \textbf{计算复杂度}基于流形的方法往往涉及复杂计算，难以扩展到大规模数据集。
  \item \textbf{模型假设}准确描述流形结构并不容易，错误的假设可能导致错误的推断。
  \item \textbf{维数估计}数据的真实内在维数难以确定，估计结果会影响基于流形的方法的效果。
\end{itemize}
未来研究可着眼于开发效率更高的算法、提高方法对模型误设的鲁棒性，以及建立准确估计维数的方法。

\section{流形上的样本统计量}

基本的样本统计量及其函数变换是欧氏空间统计学的主要内容。要建立流形上的统计理论，首先需要研究流形上的样本统计量，例如卡尔歇均值（Karcher mean）和弗雷歇均值。

在经典欧氏统计中，给定数据 $x_1,x_2,\ldots,x_n\in\R^d$，样本均值定义为
\[
  \bar x=\frac{1}{n}\sum_{i=1}^{n}x_i,
\]
它是欧氏距离平方和的最小化点，即
\[
  \bar x=\operatorname*{arg\,min}_{\mu\in\R^d}
           \sum_{i=1}^{n}\lVert x_i-\mu\rVert^2.
\]

% PDF p. 274
当数据位于流形 $M$ 上时，例如球面 $S^2$ 上的点或形状空间中的点，由于 $M$ 没有全局线性结构，算术平均值无从定义。这时用 $M$ 的内在几何来定义均值及其他统计量。

\subsection{弗雷歇均值（内在均值）与卡尔歇均值}

\subsubsection{弗雷歇均值的定义}

设 $(M,d)$ 是完备度量空间，其中 $d$ 是黎曼流形 $M$ 上的测地距离。给定样本点 $x_1,x_2,\ldots,x_n\in M$，定义弗雷歇函数（Fréchet function）
\[
  F(p)=\frac{1}{n}\sum_{i=1}^{n}d^2(p,x_i),\qquad p\in M.
\]
弗雷歇均值（或内在均值，intrinsic mean）定义为
\[
  \hat\mu=\operatorname*{arg\,min}_{p\in M} F(p).
\]
当最小化点唯一时，有时也称其为卡尔歇均值。有些文献只在无法保证全局凸性的情况下，将对 $F$ 执行梯度下降所得的局部唯一最小化点称为卡尔歇均值。

\subsubsection{直观解释与示例}

直观上，弗雷歇均值是流形上“最能代表”数据的点，它使测地距离平方和最小，正如算术平均值使 $\R^d$ 中的欧氏距离平方和最小。

例如，若数据位于球面 $S^2$ 上，弗雷歇均值就是球面上使到各数据点的大圆距离平方和最小的点。若数据分布得不太分散，即落在一个测地凸（geodesically convex）区域内，弗雷歇均值唯一。弗雷歇均值推广了欧氏空间中的算术平均值。

\subsection{流形上的经验分布函数}

\[
  P_n=\frac{1}{n}\sum_{i=1}^{n}\delta_{x_i},
\]
其中 $\delta_{x_i}$ 是点 $x_i$ 处的狄拉克测度（Dirac measure）。这一定义可直接推广到流形上。即给定 $x_1,x_2,\ldots,x_n\in M$，$M$ 上的经验测度（empirical measure）为
\[
  P_n=\frac{1}{n}\sum_{i=1}^{n}\delta_{x_i}.
\]
于是，弗雷歇函数可写成积分形式
\[
  F(p)=\int_M d^2(p,x)\,dP_n(x)
      =\frac{1}{n}\sum_{i=1}^{n}d^2(p,x_i),\qquad p\in M.
\]
% PDF p. 275
这一写法把 $P_n$ 视为 $M$ 上真实概率测度的估计，为建立渐近理论提供了起点，例如研究弗雷歇均值的相合性和中心极限定理。

\subsection{流形上经验分布的方差}

设 $M$ 为完备黎曼流形，测地距离为 $d(\cdot,\cdot)$。设 $x_1,x_2,\ldots,x_n\in M$ 是独立同分布的随机变量，其共同分布为概率测度 $P$。样本的经验分布为
\[
  P_n=\frac{1}{n}\sum_{i=1}^{n}\delta_{x_i},
\]
其中 $\delta_{x_i}$ 是点 $x_i$ 处的狄拉克测度。

定义 $P_n$ 的方差，一个自然的办法是利用样本弗雷歇函数
\[
  F_n(p)=\int_M d^2(p,x)\,dP_n(x)
        =\frac{1}{n}\sum_{i=1}^{n} d^2(p,X_i).
\]
我们将其最小化点称为样本弗雷歇均值（或内在均值）
\[
  \hat\mu_n=\operatorname*{arg\,min}_{p\in M}F_n(p).
\]
由此定义样本方差
\[
  \Var(P_n)=\min_{p\in M}\frac{1}{n}\sum_{i=1}^{n}d^2(p,X_i)
           =\frac{1}{n}\sum_{i=1}^{n}d^2(\hat\mu_n,X_i).
\]
换言之，方差等于各样本点到样本弗雷歇均值的测地距离平方的平均值。

\subsection{弗雷歇均值的中心极限定理}

\begin{enumerate}[label=\alph*.]
  \item \textbf{总体设定}\par
  设 $\mu$ 是分布 $\mathbb P$ 的总体弗雷歇均值，即
  \[
    \mu=\operatorname*{arg\,min}_{p\in M}\int_M d^2(p,x)\,d\mathbb P(x).
  \]
  我们假定 $\mu$ 唯一，且位于 $M$ 的一个测地凸区域内，因此黎曼对数映射在其附近有良好的定义。
  \item \textbf{中心极限定理（central limit theorem，CLT）的陈述}\par
  在欧氏空间的经典统计中，中心极限定理是概率论中最重要的结果之一。熟知的欧氏中心极限定理说明，大量随机变量的平均值近似服从正态分布，即
% PDF p. 276
  \[
    \sqrt n\,(\bar x_n-\mu)\sim N(0,\Sigma).
  \]
  欧氏中心极限定理可以推广到黎曼流形。在正则性条件下，尤其是弗雷歇函数在 $\mu$ 处的海森矩阵非奇异、且分布适当集中时，样本弗雷歇均值 $\hat\mu_n$ 在切空间 $T_\mu M$ 中满足中心极限定理。定义
  \[
    \hat\xi_n=\log_\mu(\hat\mu_n)\in T_\mu M,
  \]
  其中 $\log_\mu$ 为黎曼对数映射（附录 6A）。那么
  \[
    \sqrt n\,\hat\xi_n=\sqrt n\,\log_\mu(\hat\mu_n)
  \]
  依分布收敛到 $T_\mu M$ 中均值为零的多元正态分布。以符号表示为
  \[
    \sqrt n\,\log_\mu(\hat\mu_n)\xrightarrow{d}N(0,\Sigma),
  \]
  其中 $\Sigma$ 是 $T_\mu M$ 上的协方差算子（covariance operator）。可以将 $\Sigma$ 看成同时取决于数据分布和由 $M$ 的几何结构导出的曲率项（附录 6B）。\translatornote{原文此处引用附录 6B。附录 6B 讨论黎曼对数映射，附录 6C 给出切空间中 $\Sigma$ 的定义。}
\end{enumerate}

\section{黎曼流形上数据的双样本检验}

检验两组样本是否来自同一分布是统计学中的经典问题。双样本检验在欧氏数据和流形数据上都有广泛应用（Chu and Da 2024; Cheng and Xie 2024; Wang et al.\ 2022; Daniel et al.\ 2013; Wynne and Duncan 2022）。在很多应用中，例如异常检测、变点检测、遗传数据、组学数据和单细胞数据的差异分析、扩散张量成像、形状分析或方向统计，观测值位于弯曲空间中。人们已为黎曼流形上的双样本检验提出若干统计量。本章着重讨论切空间双样本检验和基于能量距离的双样本检验。

\subsection{切空间双样本检验（流形上的霍特林 \texorpdfstring{$T^2$}{T²} 检验）}

若假定数据比较集中，使流形在局部看起来近似欧氏空间，一种常用策略是将数据投影到切空间，再使用经典多元检验，例如霍特林 $T^2$ 检验。步骤如下。

\begin{enumerate}
  \item \textbf{映射到切空间}\par
  给定黎曼流形 $M$ 上的两组样本
  \[
    X=\{x_1,x_2,\ldots,x_{n_1}\}\quad\text{和}\quad
    \{y_1,y_2,\ldots,y_{n_2}\},
  \]
% PDF p. 277
  选择一个参照点 $\mu$，通常取合并样本的弗雷歇均值。然后用对数映射将每个观测值映射到切空间 $T_\mu M$，
  \[
    u_i=\log_\mu(x_i),\qquad v_j=\log_\mu(y_j).
  \]

  \item \textbf{计算样本均值与协方差}\par
  首先计算切空间中的样本均值，
  \[
    \bar u=\frac{1}{n_1}\sum_{i=1}^{n_1}u_i,\qquad
    \bar v=\frac{1}{n_2}\sum_{j=1}^{n_2}v_j.
  \]
  定义合并协方差矩阵
  \[
    S=\frac{\sum_{i=1}^{n_1}(u_i-\bar u)^2+
              \sum_{j=1}^{n_2}(v_j-\bar v)^2}{n_1+n_2-2}.
  \]
  \translatornote{若把切向量写作列向量，此处向量的平方应理解为外积，例如 $(u_i-\bar u)(u_i-\bar u)^{\mathsf T}$。}

  \item \textbf{检验统计量}\par
  经典的霍特林 $T^2$ 统计量定义为
  \[
    T^2=\frac{n_1n_2}{n_1+n_2}(\bar u-\bar v)^{\mathsf T}S^{-1}(\bar u-\bar v).
  \]

  \item \textbf{零假设下的分布}\par
  如果数据在切空间中的表示近似服从多元正态分布，且流形曲率相对于数据的离散程度较小，那么在零假设下（即两组样本来自同一分布），有
  \[
    \frac{n_1+n_2-d-1}{(n_1+n_2-2)d}\sim F_{d,n_1+n_2-d-1},
  \]
  其中 $d=\dim(T_\mu M)$ 为切空间维数，与流形的维数相同。\translatornote{原式左端缺少 $T^2$。经典的 $F$ 分布换算式是用左端比例因子乘以 $T^2$。}

  \item \textbf{模拟}\par
  下面是一项模拟研究。我们用集中程度较低的冯·米塞斯–费希尔（von Mises–Fisher，vMF）模型在球面 $S^2$ 上生成数据，使数据较为分散。两组具有不同的弗雷歇均值，按球面内在距离衡量时，它们有明显差别。由于数据分散，$\R^3$ 中的欧氏距离（弦长）相对于内在距离（测地距离）有所失真。在这种情况下，预期考虑流形几何的检验有更高的检验功效。

  模拟采用以下设置。\par
  \textbf{第 1 组}以 $\mu_1=(0,0,1)$ 为中心。\par
  \textbf{第 2 组}以 $\mu_2=(0.5,0,0.866)$ 为中心。
% PDF p. 278
  （注意 $\mu_2$ 位于 $S^2$ 上，它的坐标来自北极旋转 $30^\circ$。）

  \textbf{集中参数}取 $\kappa=2$，数值较小，因而数据相当分散。

  然后计算两种霍特林 $T^2$ 检验。\par
  \textbf{欧氏霍特林 $T^2$ 检验}直接用原始三维坐标计算，忽略球面几何。\par
  \textbf{流形霍特林 $T^2$ 检验}先计算 $S^2$ 上合并样本的弗雷歇均值，用对数映射把样本点映到切空间，再在该切空间中进行霍特林 $T^2$ 检验。

  一次典型模拟得到如下结果（不同次模拟的具体数值会变化）。\par
  \textbf{$S^2$ 上的欧氏霍特林 $T^2$ 检验} $T^2\approx 12.34$，$P$ 值约为 $0.084$。\par
  \textbf{$S^2$ 上的流形霍特林 $T^2$ 检验} $T^2\approx 18.56$，$P$ 值约为 $0.015$。

  因此，在本例中，流形检验得到的 $P$ 值比欧氏检验小得多，表明它对两组之间的内在差异更敏感，即具有更高的检验功效。
\end{enumerate}

\subsection{流形上基于能量距离的双样本检验}

在局部高斯近似不适用，或者希望采用完全非参数的方法时，可以采用基于距离的检验。一种常用方法是能量距离（energy distance）检验，采用测地距离将其推广到流形。适用于非欧氏数据和高维数据的现有非参数双样本检验通常借助某种相似性度量，常用的是欧氏距离。常见的非参数方法包括基于点间距离的检验统计量、基于核的最大均值差异（maximum mean discrepancy，MMD）检验，以及基于图的双样本检验（Chu and Da 2024）。

下面介绍基于能量距离的流形双样本检验。这是一种适用于一般数据对象的非参数双样本检验，它考虑数据的内在几何。用测地距离作为由数据决定的度量，可以在尊重流形结构的同时刻画点间距离。检验统计量的表现类似于测地距离度量。

\begin{enumerate}
  \item \textbf{测地距离}\par
  以 $d(x,y)$ 表示 $M$ 上 $x$ 与 $y$ 之间的测地距离。

  \item \textbf{能量距离统计量}\par
  定义统计量
  \[
    \begin{aligned}
      \varepsilon(X,Y)
        &=\frac{2}{n_1 n_2}\sum_{i=1}^{n_1}\sum_{j=1}^{n_2}d(x_i,y_j)
         -\frac{1}{n_1^2}\sum_{i=1}^{n_1}\sum_{i'=1}^{n_1}d(x_i,x_{i'})\\
        &\quad-\frac{1}{n_2^2}\sum_{j=1}^{n_2}\sum_{j'=1}^{n_2}d(y_j,y_{j'}).
    \end{aligned}
  \]
  在两组样本来自同一分布的零假设下，预期能量距离 $\varepsilon(X,Y)$ 接近零。

  \item \textbf{分布与推断}\par
  实际中 $\varepsilon(X,Y)$ 的零分布较复杂，因此常用置换检验。
% PDF p. 279

  \textbf{置换步骤}将合并数据随机重新分到样本量为 $n_1,n_2$ 的两组中，对每次置换计算 $\varepsilon$，再将观测统计量与所得经验分布比较。

  \textbf{渐近行为}在适当的正则性条件下，检验统计量（或经适当缩放后的统计量）依分布收敛到若干卡方随机变量的加权和。不过，置换方法更为灵活，对分布的假设也更少，因此应用较广。
\end{enumerate}

\begin{statement}{例 6.1}{}
设数据（$n_1=n_2=100$）在球面 $S^2$ 上由以下模型生成。

\textbf{第 1 组}以 $\mu_1=(0,0,1)$ 为中心。

\textbf{第 2 组}以 $\mu_2=(0.5,0,0.866)$ 为中心。

（注意 $\mu_2$ 位于 $S^2$ 上，它的坐标来自北极旋转 $30^\circ$。）

\textbf{集中参数}取 $\kappa=2$，数值较小，因而数据相当分散。

我们重复模拟 500 次。每次用置换检验计算基于能量距离的双样本检验的 $p$ 值，以检验两组之间的差异。

500 次重复完成后，计算平均 $p$ 值和拒绝率（$p<0.05$ 的模拟所占比例），汇总结果。

例如，一次典型模拟可能得到以下结果。
\begin{itemize}
  \item \textbf{平均 $p$ 值}约为 $0.012$。
  \item \textbf{$\alpha=0.05$ 时的拒绝率}约为 $0.88$。
\end{itemize}
\end{statement}
